\section{A Candidate Transition Operator for Mythic Narratives}
\label{sec:transition-operator}

\subsection{Epistemic scope}

The purpose of this section is not to prove a theological proposition by
numerical coincidence. We ask a narrower question: whether recurrent
narrative motifs admit the same finite ordered structure after their
culture-specific labels are removed. Numerical coincidence alone is
therefore explicitly insufficient.

Let a mythic motif be represented by a finite ordered event structure
\[
\mathcal{M}=(V,\prec,E,\lambda),
\]
where \(V\) is a finite state set, \(\prec\) is narrative order, \(E\) is the
set of directed transitions, and \(\lambda:E\to\mathcal{T}\) assigns a typed
transition such as \emph{withdrawal}, \emph{liminality}, \emph{emergence},
\emph{cessation}, or \emph{renewal}.

Two motifs are called structurally isomorphic only when a bijection
\(f:V_1\to V_2\) preserves both order and transition type:
\[
u\prec_1 v \iff f(u)\prec_2 f(v),
\qquad
\lambda_1(u,v)=\lambda_2(f(u),f(v)).
\]
If transition order is preserved only after collapsing or inserting
intermediate states, we use the weaker term \emph{order-homomorphism}.
Shared semantic motifs without such a map are classified merely as
\emph{analogy}. A repeated integer without preserved transition structure is
classified as a \emph{number match only}.

\subsection{The \texorpdfstring{\(123\mid45\mid6\to7\)}{123|45|6 to 7} partition}

The working Genesis-derived candidate is
\[
123\mid45\mid6\longrightarrow7.
\]
Its active block lengths are
\[
\mathbf b=(3,2,1,0),
\]
where the final zero denotes that the seventh position is not being counted
as an additional active creation block. Hence
\[
3+2+1=6=T_3,
\qquad
3>2>1>0.
\]
We define the candidate global horizon
\[
H:=6
\]
and the post-closure state
\[
S:=7.
\]
This notation is a model imposed for testing. Genesis 2:2--3 directly
supports only the weaker statement that completion is associated with six
and the seventh is distinguished by cessation/rest. It does \emph{not}
itself attest the \(3\mid2\mid1\) partition.

\subsection{Third-day transition operator}

Hosea 6:2 and 1 Corinthians 15:4 independently place a raising/revival event
at a ``third day.'' This motivates, but does not prove, a candidate local
threshold operator
\[
\Theta_3:
X_{\mathrm{liminal}}
\longrightarrow
X_{\mathrm{renewed}}.
\]
The empirical question is not whether the integer \(3\) occurs frequently,
but whether independent narratives preserve the same ordered transition
signature around that marker.

\subsection{Six-to-seven regime change}

Philo of Alexandria gives unusually explicit ancient support for treating
the six/seven distinction structurally rather than as a simple clock
duration. In \emph{Allegorical Interpretation} I.2--5 he rejects a naive
six-day temporal reading, associates time with heavenly motion, treats six
as a perfect/order number, and assigns seven a qualitatively different
role. This is evidence for an ancient \emph{regime distinction}; it is not
evidence that the intended clock was a galactic orbit.

Accordingly we define only the candidate transition
\[
\mathcal C_6:
\mathrm{ACTIVE}
\longrightarrow
\mathrm{CLOSED},
\qquad
\mathcal S_7:
\mathrm{CLOSED}
\longrightarrow
\mathrm{POST\text{-}CLOSURE}.
\]

\subsection{Cycle-defined time}

For a phase variable \(\phi_X(t)\), a cycle-defined ``day'' may be written
abstractly as
\[
D_X=\inf\left\{
t>0:
\phi_X(t)-\phi_X(0)=0 \pmod{2\pi}
\right\}.
\]
For uniform angular velocity \(\omega_X\neq0\),
\[
D_X=\frac{2\pi}{|\omega_X|}.
\]
This equation is a generic mathematical definition of recurrence. It does
not by itself identify any ancient ``divine day'' with a particular modern
astronomical period.

\subsection{Amaterasu as a provenance stress test}

The Amaterasu cave narrative is especially useful because it can expose a
methodological error. The old Japanese narrative securely contains the
withdrawal/darkness/emergence motif. The specific three-days-and-three-nights
duration used in some comparisons belongs to a later blind-\emph{biwa}
ritual tradition, in which three days and nights of \emph{mikagura} occur
before the cave still remains closed and further action precedes emergence.

Therefore
\[
\text{canonical cave motif}
\not\Rightarrow
\text{canonical three-day duration}.
\]
Any implementation that silently promotes the later duration into the
oldest textual layer fails the provenance test.

\subsection{GREMLIN routing}

The candidate search was routed through the live GREMLIN/PhaseNav layer.
The first two candidate species were:
\begin{enumerate}
  \item \textbf{OWL}: primary-evidence methodology;
  \item \textbf{SPIDER}: structural relation/dependency.
\end{enumerate}
The ordering is adopted as a methodological constraint:
\[
\mathrm{SOURCE}
\rightarrow
\mathrm{PROVENANCE}
\rightarrow
\mathrm{EVENT\ GRAPH}
\rightarrow
\mathrm{RELATION\ CLASS}.
\]
GREMLIN output remains candidate-only; it is not an evidential authority.

\subsection{Falsification criteria}

The \(3\to2\to1\to0\) hypothesis is weakened or rejected when:
\begin{enumerate}
  \item the partition must be changed separately for each tradition;
  \item only the same numeral survives while transition graphs disagree;
  \item a later textual layer must be promoted to an older one to obtain a
        match;
  \item the proposed isomorphism fails preservation of order or transition
        type;
  \item astronomical parameters must be tuned after observing target dates.
\end{enumerate}

The accompanying executable tests encode these constraints so that a
narrative resemblance cannot be promoted to an isomorphism merely because
the same integer occurs in both sources.
